\section{Monadicity of \texorpdfstring{\cstars{}s}{σC*-algebras} over \texorpdfstring{\cstar{}s}{C*-algebras}}\label{sec:monadic}
	This section proves the central result of the paper: the forgetful functor $$U \colon \SigmaC \to \Calg$$ is monadic (\cref{thm:U_monadic}), and the associated monad constructs the \emph{Pedersen--Baire envelope} of a \cstar.
	We then record the two consequences used repeatedly below: $\SigmaC$ is complete with limits created by $U$ (\cref{cor:complete}), and $T$ preserves reflexive coequalizers, whence $\SigmaC$ is cocomplete (\cref{prop:cocomplete}).

We recall the following adjunction, first shown in \cite{wright76regular}. The present formulation is taken from \cite[Equation (3.3.4)]{fritz2025categories}.
\begin{theorem}\label{thm:adj_sigma}
	There is an adjunction
	\begin{equation*}
		% \begin{tikzcd}
		%  \Calg \ar[rr,phantom, "\bot"]\ar[rr,bend left=20,"F"] && \SigmaC \ar[ll,,"U",bend left=20]
		% \end{tikzcd}
		\adPair{\Calg}{\SigmaC}{F}{U}
	\end{equation*}
	where $U$ forgets $\sigma$-normality, and $F$ sends a \cstar{} $\A$ to its \emph{Pedersen--Baire envelope}.
\end{theorem}
Explicitly, the \newterm{Pedersen--Baire envelope} $\A^{\infty}$ of a \cstar{} $\A$ is the $\sigma$-closure of $\A$ inside its enveloping \wstar{} $\A^{**}$~\cite[Definition 3.3.1]{fritz2025categories} or, equivalently, the $\sigma$-closure of $\A$ in $B(H_u)$ for the universal representation $(H_u,\pi_u)$~\cite[Definition 2.2.23]{saito15monotone}. Diagrammatically,
\[
	\xymatrix{
		\A \, \ar@{^{(}->}[r] & \A^{\infty}\,\ar@{^{(}->}[r]^{(\star)}& \A^{**}},
\]
where the right-hand side inclusion $(\star)$ is $\sigma$-normal.
The content of the adjunction of \cref{thm:adj_sigma} is exactly the universal property of $\A^{\infty}$: every $*$-homomorphism from $\A$ into a \cstars{} extends uniquely, and $\sigma$-normally, along $\A \hookrightarrow \A^{\infty}$. While the notation $\A^{\infty}$ is already used in the literature, in this paper we instead opt for $UF\A$ and $F\A$, mindful of the associated adjunction.

\begin{notation}
	Below, $F$ and $U$ always mean the two functors of the previous statement; the associated monad $UF$ will be denoted by $T$ for simplicity; $\epsilon \colon FU \to \id_{\SigmaC}$ and $\eta\colon \id_{\Calg} \to UF$ will denote the counit and unit, respectively, and $\mu\coloneqq U\epsilon F$ is the monad multiplication.
\end{notation}

We now wish to show that $U$ is \emph{monadic}. To say that $U$ is monadic is to say that a \cstars{} is \emph{nothing more than} a \cstar{} equipped with a coherent choice of the countable suprema forced by the Pedersen--Baire envelope: the additional structure is algebraic over $\Calg$, carried by operations and equations rather than by a property. This is what buys, in the sections that follow, limits created by $U$, free objects, a tensor product of algebras (\cref{sec:tensor}) and a presentability theorem (\cref{sec:accessible}), none of which is visible from the definition of a \cstars{}.

	A warning on terminology may spare the reader a search: a substantial part of the literature on monadicity, including \cite{barr2005toposes} and \cite{linton1969coequalizers} as cited below, says \emph{triple} and \emph{tripleable} where we say monad and monadic; \cite{barr2005toposes} does not use the word ``monad'' at all.
	
Monadicity of $U$ can be argued directly, showing that the unique comparison functor $K\colon \SigmaC \to \EM(T)$ deduced from \cite[p. 222]{macdonald2003aspects} is fully faithful and essentially surjective (the two halves of the proof of \cref{thm:U_monadic} below), thus an equivalence of categories.
 The functor $K$ is the only one filling the commutative diagram,
\[
	\xymatrix{
		\SigmaC \ar[rr]^{K} \ar[dr]_{U} && \EM(T) \ar[dl] \\
		& \Calg
	}
\]
where the unlabeled arrow is a forgetful functor.
Explicitly, $K$ sends a \cstars{} $\A$ to the carrier given by $U\A$ equipped with the algebra structure given by the counit $U\epsilon_{\A} \colon TU\A \to U\A$.
\begin{theorem}\label{thm:U_monadic}
	$U$ is monadic.
\end{theorem}
\begin{proof}
	We divide the proof in two steps.
	\setlist[description]{font=\normalfont\itshape}
	\begin{description}
		\item[{$K$ is fully faithful}]   By a proposition due to Beck (see, e.g., \cite[p.~226]{macdonald2003aspects}), $K$ is fully faithful if and only if $\epsilon_{\B}$ is a regular epimorphism.
		      We seek to show that this is always the case.

		      First of all, we recall that $\epsilon_{\B}\colon FU\B \to \B$ is a surjective map in $\SigmaC$.
		      This was proven in \cite[Theorem 5.4.5]{saito15monotone} (see \cite[Theorem 3.3.6]{fritz2025categories} for the explicit statement).
		      As is standard for algebras, every surjective map in fact identifies its image with the quotient via its kernel, i.e. $\B \cong FU\B / \ker(\epsilon_{\B})$.

		      We are therefore left to show that this allows one to write $\epsilon_{\B}$ as a coequalizer. To this end, we consider two cases. If $\ker(\epsilon_{\B})$ does not have a unit, we consider the unitization $\ker(\epsilon_{\B})_1$ of $\ker(\epsilon_{\B})$, which is a \cstars{} as shown in \cite[Proposition 2.2.10]{saito15monotone}.
		      There are two obvious maps $\ker(\epsilon_{\B})_1 \to FU\B$.
		      The first one, which we call $i$, sends $\lambda 1+x$, where $\lambda \in \bbC$ and $x\in \ker(\epsilon_{\B})$, to itself via the inclusion $\ker(\epsilon_{\B}) \subset FU\B $, since $FU\B$ is unital. The second one, denoted by $\ell$, forgets $x$: in other words, it sends $\lambda 1 +x$ to $\lambda 1$. This is well-defined as $\ker(\epsilon_{\B})$ is a (maximal) ideal of its unitization $\ker(\epsilon_{\B})_1$.
		      It is now a simple check that $\epsilon_{\B}$ is the coequalizer of
		      \[
			      \xymatrix{
				      \ker(\epsilon_{\B})_1
				      \ar@<.7ex>[r]^{i}
				      \ar@<-.7ex>[r]_{\ell}
				      & FU\B
			      }
		      \]
		      since the image of their difference $i-\ell$ exactly corresponds to $\ker(\epsilon_{\B})$.
		      If instead $\ker(\epsilon_{\B})$ has a unit, we consider $\ker(\epsilon_{\B})\oplus \bbC$. In this case, we take $i\colon \ker(\epsilon_{\B})\oplus \bbC \to FU\B$ sending $(x,\lambda)$ to $x+\lambda (1-e)$, where $e$ denotes the unit of $\ker(\epsilon_{\B})$, and $\ell$ the composition $\ker(\epsilon_{\B})\oplus \bbC \to \bbC\to FU\B$, where the first map simply sends $(x,\lambda)$ to $\lambda$.
		\item[{$K$ is essentially surjective}] Let us consider an algebra in $\EM(T)$. By definition, this is given by a \cstar{} $\A$ together with a morphism $\phi\colon T\A\to \A$ compatible with the monad multiplication and unit.
		      Explicitly,
		      \begin{equation}\label{eq:EM_algebra}
			      \phi \eta_{\A}=\id_{\A}\qquad\text{and}\qquad\phi T(\phi)= \phi \mu_{\A},
		      \end{equation}
		      where $\eta_{\A}\colon \A \to T\A$ denotes the standard inclusion, while $\mu_{\A}=U(\epsilon_{F\A})\colon TT\A \to T\A$ is the $\sigma$-normal map such that $\mu_{\A}\eta_{T\A}=\id_{T\A}$.

		      Let us show that $\A$ satisfying \eqref{eq:EM_algebra} is a \cstars{} and moreover that $\phi$ is $\sigma$-normal.
		      In other words, $\A = U\A$ and $\phi=U\phi$.
		      Let $(a_n)$ be a norm-bounded monotone increasing sequence in $\A$.
		      Since $\eta_{\A}$ preserves the positive ordering, the image of the sequence $(\eta_{\A}(a_n))$ is a norm-bounded monotone increasing sequence in $T\A$, which is a \cstars{}.
		      We can then consider $\phi(\sup_n^{T\A}(\eta_{\A}(a_n)))\in \A$.
		      Now, given any $b\ge a_n$ for all $n$, $\eta_{\A}(b) \ge \sup_n^{T\A}(\eta_{\A} (a_n))$ and then
		      \[
			      b=\phi(\eta_{\A}(b)) \ge \textstyle\phi\left(\sup_n^{T\A}(\eta_{\A}(a_n)\right)
		      \]
		      Finally, $\sup_n^{\A} a_n $ exists and is given by $\phi (\sup_n^{T\A}(\eta_{\A} (a_n)))$.
		      To show that $\phi$ is $\sigma$-normal, we use the identity given by the algebra multiplication: consider $\sup^{TT\A}_n \eta_{T\A} (x_n)\in TT\A$ for a norm-bounded monotone increasing sequence $(x_n)\subseteq T\A$.
		      Then
			      \begin{align}
				      \textstyle\phi\left(\sup_n^{T\A} \eta_{\A} \phi(x_n)\right) & = \textstyle\phi T(\phi) \left(\sup_n^{TT\A} \eta_{T\A} (x_n)\right)  \tag{definition of $T(\phi)$}\\
				                                                        & = \textstyle\phi \mu_{\A} \left(\sup_n^{TT\A}\eta_{T\A} (x_n)\right) \tag{\ref{eq:EM_algebra}}\\
				                                                        & = \textstyle\phi \left(\sup_n^{T\A} x_n\right)                       \tag{$\sigma$-normality of $\mu$, and $\mu \circ \eta T =\id_{T\A}$}
			      \end{align}
		      Since we just established that $\sup_n^{\A} a_n = \phi (\sup_n^{T\A} (\eta_{\A}(a_n)))$, the left-hand side is equal to $\sup^{\A}_n (\phi(x_n))$.
		      Finally, $\sup^{\A}_n (\phi(x_n)) = \phi(\sup^{T\A}_n x_n)$, as wanted.

		      We now claim that $\phi$ is in fact $\epsilon_{\A}$, from which essential surjectivity of $K$ follows.
		      To this end, consider the set
		      \[M\coloneqq \lbrace a\in T\A \mid \phi(a)=\epsilon_{\A}(a)\rbrace, \]which certainly contains $\eta_{\A} (\A)$ by the equations $\phi \eta_{\A}=\id_{\A}$ and $U(\epsilon_{\A}) \eta_{\A}=\id_{\A}$.
		      Since both $\phi$ and $\epsilon_{\A}$ respect suprema, $M$ is also closed under suprema.
		      Similarly, one also proves that $M$ is a $*$-subspace, hence it is a $\sigma$-subspace.
		      As the $\sigma$-closure of $\eta_{\A} (\A)$ in $T\A$ is the whole $T\A$ (see the discussion after \cref{thm:adj_sigma}), we obtain that $M=T\A$, thus concluding the proof.\qedhere
	\end{description}
\end{proof}


By \cite[Proposition 4.3.1]{borceux94handbook2}, since $\Calg$ has all limits~\cite[Section 2]{khoshkam2002categorical}, \cref{thm:U_monadic} yields the following.

\begin{corollary}\label{cor:complete}
	$\SigmaC$ is complete, and limits are created, in the sense of \cite[\S V.1]{cwm}, by $U$.
\end{corollary}

\begin{remark}[Monadicity for cpu maps, and commutative algebras]\label{rem:monadic_cpu}
	The adjunction of \cref{thm:adj_sigma} can be extended to \cpu{} maps via the \emph{Stinespring dilation theorem}, as shown in \cite[Corollary 3.3.8]{fritz2025categories}.
	Explicitly, we have
	\[
		\adPair{\Calgcpu}{\SigmaCcpu}{\bar F}{\bar U}
	\]
	extending the one in \cref{thm:adj_sigma}.
	The reasoning of \cref{thm:U_monadic} can now be copied verbatim to ensure that the extended $\bar U $ is monadic as well: in particular, note that the proof of essential surjectivity of $K$ only requires $\phi$ to preserve the order.

	Moreover, by restricting the whole discussion to the commutative case, since $F\A$ is always commutative provided $\A$ is (\cite[Remark 5.3.4(iii)]{saito15monotone}), the same argument proves monadicity of the forgetful functors $\cSigmaC\to\cCalg $ and $\cSigmaCcpu\to \cCalgcpu $.
\end{remark}

We now seek to show that $\SigmaC$ is also cocomplete.
To this end, we show two preparatory lemmas, where we need to consider nonunital settings.
Recall from \cite[\S2.2]{saito15monotone} that a \newterm{nonunital \cstars{}} is a nonunital \cstar{}\footnote{Here, ``nonunital'' means ``not necessarily unital''; thus unital \cstar{}s are particular instances of nonunital \cstar{}s.} such that all norm-bounded monotone increasing sequences admit a supremum.
This is different from requiring that all monotone increasing sequences bounded by an element admit a supremum: this second option gives rise to \emph{pseudo monotone $\sigma$-complete (nonunital) \cstar{}s} and it allows one to consider the compact operators over a separable Hilbert space as an example, while the definition we adopt here does not (see \cite[Example 2.2.11]{saito15monotone} for details).

\begin{lemma}\label{lem:sigmaC_ideal}
	Let $\B$ be a \cstars{}, and let $I$ be a closed ideal in $\B$ and a nonunital \cstars{}.
	Then $I$ is a $\sigma$-ideal.
\end{lemma}
\begin{proof}
	Given a norm-bounded monotone increasing sequence $(a_n)\subseteq I$, we have $\sup_n^I a_n \ge \sup_n^{\B}a_n \ge a_n$ because $I\subseteq \B$ and both are \cstars{}s ($I$ nonunital). Since closed ideals are \emph{hereditary}~\cite[Theorem 3.2.1]{murphy90cstar}, that is, $0 \le b \le a$ with $a \in I$ and $b\in\B$ forces $b\in I$,
	we deduce that $\sup_n^{\B}a_n $ must also be contained in $I$, hence $\sup_n^{\B} a_n = \sup_n^I a_n$.
\end{proof}

\begin{lemma}\label{lem:reflexive_coeq}
	Consider \cstars{}s $\A$ and $\B$. Let us consider the diagram
	\[
		\xymatrix{
		\A
		\ar@<1.2ex>[r]^f
		\ar@<-1.2ex>[r]_g
		&
		\B
		\ar@{->}[l]|-{s}
		}
	\]
	where $f$ and $g$ are \starshom{}s, $s$ is any \cpu{} map, and $fs=gs=\id_{\B}$. The reflexive coequalizer $\C$ of the above diagram in $\Calgcpu$ is a \cstars{}, and moreover the induced $\B \to \C$ is a surjective \starshom{}.

	In other words, the forgetful functor $\SigmaC \to \Calgcpu$ preserves coequalizers that become reflexive after applying the functor.
	Consequently, the same holds for the forgetful functors $\SigmaC\to \Calg$ and $\SigmaC \to \SigmaCcpu$.
\end{lemma}
This statement is similar to \cite[Proposition 4.4]{furber2024monad}.
In particular, this result could be used to reprove monadicity of $U$ via Beck's theorem as stated in \cite[Theorem 3.14]{barr2005toposes}.

\begin{proof}
	The whole proof consists in proving that the $*$-subspace $I\coloneqq {\Im(f-g)}$ is a $\sigma$-ideal.
	Indeed, if this is the case, the coequalizer is given by the quotient map $\B \to \B/I$, and the required monotone properties follow immediately from \cite[Proposition 2.2.21]{saito15monotone}.

	First, we note that $I$ is a two-sided ideal.
	Indeed, given $a\in \A$ and $b\in \B$, then
	\[ b(f(a)-g(a)) = f(s(b)) f(a) - g(s(b)) g(a) = f(s(b)a)-g(s(b)a),\]
	and analogously on the other side.

	To show it is both closed and $\sigma$-closed, we consider the (nonunital) \starshom{}
	\[
		f'\colon
		\xymatrix{\ker(g)\, \ar@{^{(}->}[r] & \A \ar[r]^f& \B}
	\]
	and claim that its image is precisely $I$.
	Indeed, given $k \in \ker(g)$, $f'(k)=f(k)-g(k)\in I$.
	Conversely, if $b\in I$, there is $a\in \A$ such that $b=f(a)-g(a)$.
	Now, by selecting $a'=a-sg(a)$, the equations $g(a')=0$ and $f(a')=b$ are directly verified.
	Finally, $I=\Im(f')$.
	Consequently, $I$ is closed by \cite[Theorem 3.1.6]{murphy90cstar}.
	Moreover, using the first isomorphism theorem for \cstar{}s, $I = {\Im(f')}\cong \ker(g)/\ker(f')$, i.e.~it is a quotient of (nonunital) \cstars{}s, thus a nonunital \cstars{} by \cite[Proposition 2.2.21]{saito15monotone}.
	We have shown that $I\subseteq \B$ satisfies the hypotheses of \cref{lem:sigmaC_ideal}, so it is $\sigma$-closed.
\end{proof}

\begin{proposition}\label{prop:cocomplete}
	$T$ preserves reflexive coequalizers.
	In particular, $\SigmaC$ is cocomplete.
\end{proposition}
\begin{proof}
	Consider $f,g\colon \A\to\B$ and $s\colon \B \to \A$ with $gs=fs=\id_{\B}$ in $\Calg$.
	We call $\C$ and $\D$ the reflexive coequalizers of the pairs $(f,g)$ and $(Tf,Tg)$, respectively.
	By \cref{lem:reflexive_coeq}, $\D$ is a \cstars{} and moreover the cocone map $T\B \to \D$ is a \starshom{}.
	Precomposing with $\eta_{\B}\colon \B \to T\B$, we obtain a map $\B \to \D$ that coequalizes $f$ and $g$.
	This entails the existence of a unique map $\C \to \D$, which lifts to a map $T\C \to \D$ because $\D$ is a \cstars{}.
	Since the necessary commutativities are directly inferred, $T\C$ must also be the coequalizer.

	The second sentence follows from Linton's theorem (\cite[Corollaries 2 and 3]{linton1969coequalizers}), since $T$ is a monad on a cocomplete category (\cite[Section 2]{khoshkam2002categorical}).
\end{proof}
\begin{remark}\label{rem:two_routes_cocomplete}
		Cocompleteness of $\SigmaC$ admits a second, less elementary proof: by \cref{thm:U_accessible} the category is locally $\aleph_1$-presentable, hence cocomplete~\cite[Definition 1.17]{adamek1994locallypresentable}.
		We have given the argument above because it is independent of the accessibility machinery of \cref{sec:accessible} and because its first assertion, that $T$ preserves reflexive coequalizers, is needed later (it is the hypothesis of \cref{thm:SMC2} used in the proof of \cref{thm:tensor_shom}) and does not follow from presentability.
		No circularity arises: the proof of \cref{thm:U_accessible} rests on \cref{thm:U_monadic} and \cref{lem:technical_access}, not on \cref{prop:cocomplete}.
\end{remark}

\begin{remark}[Monadicity of (commutative) \wstar{}s]\label{rem:monadicity_W}
	Consider the category $\Walg$ of \wstar{}s and normal (unital) $*$-homomorphisms.
	The forgetful functor $\Walg \to \Calg$ has been proven to be monadic in \cite[Theorem 4.6]{furber2024monad}.
	Moreover, \cite[Proposition 4.4]{furber2024monad} also ensures that this functor preserves reflexive coequalizers.
	In particular, $\Walg$ is cocomplete by the same reasoning of the proof of \cref{prop:cocomplete}.

	Combining this result with \cref{thm:U_monadic} and \cite[Corollary 4.5.7]{borceux94handbook2}, we immediately obtain that $\Walg\to \SigmaC$ is monadic as well.

	A similar result holds by restricting to commutative algebras.
\end{remark}

We conclude with a simple remark that will be used in the proof of \cref{thm:tensor_scpu}.
\begin{remark}\label{rem:T_surjective}
	$T$ preserves surjective maps.
	Indeed, given a surjective map $\phi \colon \A \to \B$ in $\Calg$, the map $T\phi \colon T\A \to T\B$ sends $\eta_{\A}(a) $ to $\eta_{\B}\phi(a)$. 
	Since $\phi$ is surjective, $\Im (T\phi)\supseteq \eta_{\B}(\B)$. 
	As $T\phi$ is $\sigma$-normal, its image must be $\sigma$-closed, so it must coincide with the whole $T\B$ because $T\B$ is the $\sigma$-closure of $\eta_{\B}(\B)$.
\end{remark}

